\DocumentMetadata{
  pdfversion=2.0,
  pdfstandard=ua-2,
  lang=en-US,
  tagging=on,
  tagging-setup={math/setup=mathml-SE,tabsorder=structure}
}
\documentclass[11pt,letterpaper]{article}
\usepackage[margin=0.68in,includefoot]{geometry}
\usepackage{newcomputermodern}
\usepackage{amsmath,amscd,mathtools}
\usepackage{tikz,adjustbox}
\providecommand{\square}{\mdlgwhtsquare}
\usepackage[hidelinks]{hyperref}
\hypersetup{
  pdftitle={Algebra PhD Exam, May 1996},
  pdfauthor={University of Florida Department of Mathematics},
  pdfsubject={Graduate mathematics examination},
  pdfdisplaydoctitle=true,
  bookmarksopen=true
}
\setcounter{secnumdepth}{0}
\setlength{\parindent}{0pt}
\setlength{\parskip}{0.28em}
\setlength{\emergencystretch}{3em}
\setlength{\leftmargini}{2.15em}
\setlength{\leftmarginii}{2.65em}
\setlength{\leftmarginiii}{3em}
\renewcommand{\labelenumi}{\textbf{\theenumi.}}
\pagestyle{empty}
\raggedbottom
\allowdisplaybreaks
\newenvironment{examproblems}[1]{%
  \begin{enumerate}%
  \setcounter{enumi}{#1}%
  \setlength{\topsep}{0.35em}%
  \setlength{\partopsep}{0pt}%
  \setlength{\itemsep}{0.48em}%
  \setlength{\parsep}{0.18em}%
}{\end{enumerate}}
\newenvironment{examsubparts}{%
  \begin{enumerate}%
  \setlength{\topsep}{0.18em}%
  \setlength{\partopsep}{0pt}%
  \setlength{\itemsep}{0.16em}%
  \setlength{\parsep}{0.1em}%
}{\end{enumerate}}
\newenvironment{examinstructions}{%
  \par\begingroup\itshape
}{%
  \par\endgroup\vspace{0.18em}
}
\makeatletter
\renewcommand\section{\@startsection{section}{1}{0pt}%
  {-1.25ex plus -.25ex minus -.1ex}%
  {0.55ex plus .1ex}%
  {\normalfont\large\bfseries}}
\renewcommand{\@maketitle}{%
  \newpage\null\vskip -1em%
  {\footnotesize\bfseries
    \href{https://www.ufl.edu/}{University of Florida}, \href{https://math.ufl.edu/}{Department of Mathematics}\par}%
  \vskip -0.5em%
  \tagstructbegin{tag=Title}%
    {\LARGE\bfseries\href{https://gma.math.ufl.edu/exams/algebra-phd/}{Algebra PhD Exam}, May 1996\par}%
  \tagstructend%
  \vskip 0.25em%
  \tagmcbegin{artifact}%
    \hrule height 1pt%
  \tagmcend%
  \vskip 1em%
}
\makeatother
\title{Algebra PhD Exam, May 1996}
\author{}
\date{}
\begin{document}
\maketitle
\thispagestyle{empty}
\begin{examinstructions}
Answer seven out of the following ten problems. State clearly any results you need.
\end{examinstructions}
\begin{examproblems}{0}
\item
Let~\(R\) be an integral domain, with fraction field~\(K\). Recall that an~\(R\)-submodule~\(I\) of~\(K\) is said to be invertible if \(II^{-1}=R\), where \(I^{-1}=\{x\in K\mid ix\in R\text{ for all }i\in I\}\). Show that if~\(I\) is invertible, then it is finitely generated and projective.
\item
This problem has two parts.
\begin{examsubparts}
\item[\textnormal{(a)}]
Let~\(k\) be an algebraically closed field and~\(R\) a commutative, finitely generated~\(k\)-algebra. Show that the Jacobson radical of~\(R\) is the same as the nilradical. (Hint: Nullstellensatz).
\item[\textnormal{(b)}]
Give an example of a commutative ring with identity whose Jacobson radical is not equal to its nilradical.
\end{examsubparts}
\item
Compute the Galois group of the extension \(K/\mathbb{Q}\), where~\(K\) is the splitting field over~\(\mathbb{Q}\) of the polynomial \(X^4-3\).
\item
Let~\(K\) be a field, and~\(R\) a simple, finite-dimensional~\(K\)-algebra with identity (not necessarily commutative). Show that if the identity is the only (nonzero) idempotent element of~\(R\), then~\(R\) is a division ring.
\item
Let~\(R\) be a discrete valuation ring with fraction field~\(K\). Show that~\(R\) is integrally closed in~\(K\).
\item
This problem has two parts.
\begin{examsubparts}
\item[\textnormal{(a)}]
Suppose~\(G\) is a finite~\(p\)-group and~\(S\) is a finite set on which~\(G\) acts. If \(S^G\) is the set of elements of~\(S\) fixed by~\(G\), show that \(\lvert S\rvert\equiv\lvert S^G\rvert\pmod p\).
\item[\textnormal{(b)}]
Suppose~\(G\) is a~\(p\)-group, \(k\) is a field of characteristic~\(p\), \(V\) is a finite-dimensional~\(k\)-vector space, and \(\rho:G\to\operatorname{GL}(V)\) is a homomorphism. Show that there is a nonzero vector \(v\in V\) such that \(\rho(g)(v)=v\) for all \(g\in G\). Hint: reduce to the case where~\(k\) is finite, and apply (a).
\end{examsubparts}
\item
Let~\(k\) be a field and \(G=\operatorname{GL}(3,k)\). Describe the conjugacy classes of elements of order 3 in~\(G\) if
\begin{examsubparts}
\item[\textnormal{(a)}]
\(k=\mathbb{C}\),
\item[\textnormal{(b)}]
\(k=\mathbb{R}\),
\item[\textnormal{(c)}]
\(k=\mathbb{F}_3\).
\end{examsubparts}
\item
Let~\(R\) be a commutative ring. Define what it means for an~\(R\)-module to be injective. Show that an~\(R\)-module~\(M\) is injective if and only if for every ideal \(I\subseteq R\) and~\(R\)-linear map \(f:I\to M\), \(f\) extends to a map \(R\to M\).
\item
This problem has two parts.
\begin{examsubparts}
\item[\textnormal{(a)}]
Show that there exist solvable groups of arbitrarily large derived length.
\item[\textnormal{(b)}]
Show that if~\(S\) is an infinite set, then there is no “free solvable group generated by~\(S\)”; i.e. there is no solvable group~\(F\) and map \(\alpha:S\to F\) such that composition with~\(\alpha\) induces, for any solvable group~\(G\), a bijection between the set of group homomorphisms \(F\to G\) and the set of maps \(S\to G\).
\end{examsubparts}
\item
State and prove the Hilbert Basis Theorem.
\end{examproblems}
\end{document}
